% Transcription — fonds Alexandre Grothendieck, shelfmark 101, batch 2 (pages 21-40).
% Established from the facsimile archives/batches/101/batch-02.pdf (a mirror of
% https://grothendieck.umontpellier.fr/101.pdf), per the skill
% transcribe-grothendieck. The batch file carries no cover sheet: PDF page k
% is page 20+k of the folder; the archivists' pencil numbers bottom-left
% agree wherever legible (« 34 » ... « 40 »).
%
% Pass: Opus 5.5 (claude-opus-5-5), 2026-09-24 — first pass, unchecked against the pages by a human.
%
% The folder sits in the inventory group "69-89" « Géométrie et topologie
% combinatoire », whose subject does not match this folder's (as for folder
% 100); \dating{} carries the folder's own date verbatim, and nothing is
% concluded from the group.
%
% All twenty pages are transcribed; nothing is skipped. Every leaf is his
% manuscript (blue ink, some later black ink and pencil); no one else's text
% is on any of them, so nothing is summarised.
%
% What the batch is. One continuous run on local cohomology and local
% duality over a noetherian local ring A of dimension n, taken up
% mid-argument from batch 1 (page 21 cites « Th. 1 », « Th. 3 », « cor. 1 »,
% and « § 1, prop. 4, cor. 1 » on page 39). H^i(M) is local cohomology,
% D(-) the Matlis dual, I a dualising module, Ω = Ω_A a « module
% fondamental » (D(Ω) ≃ H^n(A): the canonical module).
%   21     End of a proof: Ω Cohen-Macaulay of dim n, H^n(Ω) dualising,
%          Hom(Ω,Ω) ≃ A, Ω faithful; Cor. 4 (Ω of injective dim n); a
%          marginal Cor. 5 (codim M ≤ k ⟺ Ext^i(M,Ω) = 0 for i ≥ n-k).
%   22-25  The pairing 𝔇: H^n(A) ⊗ Ω = H^n(Ω) → I and the isomorphisms
%          𝔇_1, 𝔇_2, 𝔇_{1,M}, 𝔇_{2,M}; Prop. 9 (Ann Ω = the intersection
%          𝔑 of the primary components of « corang » n; Ω fundamental for
%          A/𝔑); Cor. 1-4 (𝔇 surjective / injective / bijective; Cor. 4:
%          bijective when A is Cohen-Macaulay or normal); a Question, with a
%          counterexample sketched (dim 2, unibranch) where Hom(Ω,Ω) is the
%          integral closure; Prop. 10 (Ω/xΩ fundamental for A/xA when x is a
%          non-zero-divisor with xH^{n-1}(A) = H^{n-1}(A)) and its Cor.
%   26-28  A a quotient of a Cohen-Macaulay (then regular) A': Ω_A
%          ≃ Ext^{n'-n}_{A'}(A, Ω_{A'}), localisation, Prop. 11 (quotients of
%          regular rings have Ω, and it localises), the Ω_A^{(i)} of (13),
%          vanishing for i ≥ 1 ⟺ A Cohen-Macaulay.
%   29     Lemma 6 (the change-of-rings spectral sequence Ext_{A'}(N,M) ⇐
%          Ext_A(N, Ext_{A'}(A,M))) and Prop. 12 (a spectral sequence
%          D(Ext^p_A(M, Ω_A^{(q)})) ⇒ H(M)).
%   30-34  His heading « 6. Caractérisation des modules fondamentaux
%          d'anneaux Coh. Mac. Anneaux de Gorenstein. »: Prop. 13 (a crossed-
%          out first version, then the version with Ext^{n-1}(k,M) = 0),
%          Cor. 1-3, Prop. 14, Prop. 15 (A fundamental over itself ⟺ H^n(A)
%          dualising ⟺ injective ⟺ comonogenic and equidimensional), a
%          struck Déf. 2 and Th. 5, Prop. 16, then « Déf 2 Anneau de
%          Gorenstein = Anneau Coh. Mac. A tel que A fond. », Prop. 17
%          (localisation, quotient by part of a system of parameters).
%   35-37  A fair-copy recapitulation: the data fixed, conditions (i)-(v) on
%          M, their proof strategy, Problems A, B, C.
%   38-39  A Theorem (under M Cohen-Macaulay the conditions are equivalent to
%          A Cohen-Macaulay and M ≃ Ω_A) and its proof by induction through a
%          non-zero-divisor x and Ext^i_A(N,M) ≃ Ext^{i-1}_{A/x}(N, M/xM).
%   40     A second recapitulation, heavily reworked in three inks; batch 3
%          (pages 41-42) continues it.
%
% Hand. Fast on 21-34 and 38-40: statements and formulas carry, connective
% prose often does not and is left \ill{}. Pages 35-37 are written fair and
% read nearly whole. Page 40 is a palimpsest: its body is given, its oblique
% marginal notes only in fragments. Page 30 is crossed through (recorded
% once there); page 33's lower half is struck line by line.
%
% Notation fixed for the batch. Ext, Hom without subscript where he writes
% none; H^i(M) as he writes it; D(-) as D; his gothic D as \mathfrak{D}; his
% gothic capital for the annihilator ideal (22-25) and the ideal of
% definition (35, 40) as \mathfrak{N} — the same glyph as on page 42, where
% batch 3 now reads it \mathfrak{N} too, for one notation across the folder; the small
% round sign for the maximal ideal as \mathfrak{m}; lengths « lg_p » as
% \mathrm{lg}_{\mathfrak{p}}; his underlined k (page 40) and x as
% \underline{} in math only; \hat for his hats. « corang » is an uncertain
% reading of a short word before « n » (22-25, 35, 38), kept uniform and
% flagged once where it first occurs. The abbreviations « Coh. Mac. »,
% « C.-M. », « fond. », « comonogène », « équidim. » are kept.
%
% Cross-references. Page 21 cites Th. 1, Th. 3 and « le cor. 1 » of batch 1;
% page 34 cites « cor. 2 à prop. 13 » and « cor. 3 »; page 38's « diagramme
% des conditions » is page 35's; page 39 cites « § 1, prop. 4, cor. 1 ». No
% explicit page number is cited. Page 40 ends mid-page; batch 3 is not
% read here beyond its header.
\documentclass[11pt,a4paper]{article}
% Le préambule commun à toutes les transcriptions du fonds Grothendieck.
%
% Il tient en une idée : **l'incertitude se compose, elle ne se résout pas.**
% Une transcription qui lisse un mot douteux a détruit la seule chose pour
% laquelle on la faisait — un lecteur doit pouvoir savoir, à chaque endroit,
% ce qui était sur le feuillet et ce qui a été ajouté. Les six macros qui
% suivent sont donc l'appareil critique, et non de la décoration.
%
% Les mêmes commandes sont reconnues par `scripts/render.mjs`, qui produit la
% vue de lecture du site. Une macro ajoutée ici doit l'être aussi là-bas,
% faute de quoi le rendu échouera bruyamment — ce qui est voulu : un rendu
% silencieusement faux est pire qu'un rendu absent.

\ProvidesPackage{grothendieck}[2026/08/08 Transcriptions du fonds Grothendieck]

\usepackage{amsmath,amssymb,amsthm}
\usepackage{mathrsfs}
\usepackage{stmaryrd}
% Les diagrammes commutatifs sont partout dans ce fonds : c'est la langue
% même de la géométrie algébrique de Grothendieck, et une transcription qui
% les rendrait en prose ne transcrirait rien.
\usepackage{tikz-cd}
% Français par défaut — c'est la langue des feuillets — mais l'anglais est
% chargé aussi : la traduction et le résumé sont des documents anglais, et la
% typographie française (espaces avant les deux-points) y serait une faute.
% Ces documents basculent par \selectlanguage{english} en tête de corps.
\usepackage[english,french]{babel}
\usepackage{fontspec}
\usepackage{geometry}
\geometry{a4paper,margin=25mm,marginparwidth=28mm}
\usepackage{marginnote}
\usepackage{xcolor}
% ulem, not soul. `soul`'s \st and \ul are built on a character-by-character
% scan that XeLaTeX's font machinery breaks: the document fails with an
% undefined control sequence rather than degrading. ulem does the same two jobs
% and is Unicode-engine safe. `normalem` stops it hijacking \emph, which the
% transcriptions use for Grothendieck's own emphasis.
\usepackage[normalem]{ulem}
\usepackage{hyperref}
\hypersetup{colorlinks=true,linkcolor=black,urlcolor=black,pdfborder={0 0 0}}

% --- Métadonnées du lot -----------------------------------------------------
% Reprises telles quelles par le script de rendu, qui les affiche en tête de
% la vue de lecture. Elles fixent ce que le fichier prétend couvrir : sans
% elles, une transcription flotte sans point d'attache dans un fonds de seize
% mille feuillets.

\newcommand{\folder}[1]{\gdef\grothFolder{#1}}
\newcommand{\batch}[1]{\gdef\grothBatch{#1}}
\newcommand{\leaves}[2]{\gdef\grothFirst{#1}\gdef\grothLast{#2}}
\newcommand{\foldertitle}[1]{\gdef\grothFoldertitle{#1}}
\gdef\grothFolder{?}\gdef\grothBatch{?}\gdef\grothFirst{?}\gdef\grothLast{?}\gdef\grothFoldertitle{}

% --- Le résumé --------------------------------------------------------------
%
% Il ouvre la lecture modernisée, sous le titre « Résumé », et il est composé
% en petit. Ce n'est pas un ornement : c'est le seul passage du document qui
% ne soit pas des mathématiques, et un lecteur doit pouvoir le reconnaître
% avant de l'avoir lu — puis le sauter s'il connaît déjà le sujet, ou s'y
% arrêter s'il ne le connaît pas. Le filet qui le suit marque la reprise du
% corps.

\newenvironment{resume}
  {\small\setlength{\parskip}{0.45em}}
  {\par\medskip\hrule\medskip}

% --- Mots-clés ---------------------------------------------------------------
%
% En anglais, en clôture du résumé de la lecture modernisée : le vocabulaire
% moderne sous lequel chercher ce que les feuillets construisent. Le manifeste
% du site (`npm run manifest`) les extrait pour étiqueter la cote entière —
% cette ligne est la source unique des tags, il n'y a pas de fichier à côté.
\newcommand{\keywords}[1]{\par\smallskip\noindent
  {\footnotesize\textsc{Keywords} --- \textit{#1}}\par}

% --- L'appareil critique ----------------------------------------------------

% Le numéro de feuillet, en marge. C'est l'ancre qui permet de régler un
% désaccord sur une formule en regardant *un* feuillet plutôt qu'une chemise
% de six cents. Le site s'en sert aussi pour tourner les pages du fac-similé
% quand on fait défiler la transcription.
\newcommand{\leaf}[1]{%
  \marginnote{\footnotesize\color{gray!70}\textsf{#1}}%
  \label{leaf:#1}\ignorespaces}

% Illisible : on ne devine pas. Un mot inventé qui se lit comme les autres est
% le pire résultat possible.
\newcommand{\ill}{\textcolor{red!60!black}{[\,\ldots\,]}}

% Lecture douteuse : le mot est donné, et signalé comme incertain.
\newcommand{\uncertain}[1]{\uline{#1}}

% Ajout de l'éditeur — une lettre manquante, un mot sous-entendu.
\newcommand{\add}[1]{\textcolor{blue!50!black}{[#1]}}

% Biffé par Grothendieck lui-même. Ce qu'il a rayé fait partie du document :
% une rature dit souvent où la pensée a changé de direction.
\newcommand{\struck}[1]{\sout{#1}}

% Note du transcripteur, sur l'état matériel du feuillet.
\newcommand{\note}[1]{\par\smallskip{\small\color{gray!60!black}\textit{[#1]}}\par\smallskip}

% Note marginale de Grothendieck, distincte du corps du texte.
\newcommand{\marginal}[1]{\par\smallskip\noindent\hspace{1em}%
  {\small\color{gray!60!black}#1}\par\smallskip}

% --- Titre ------------------------------------------------------------------

\AtBeginDocument{%
  \begin{center}
    {\large\bfseries Fonds Alexandre Grothendieck --- cote n\textsuperscript{o}~\grothFolder}\\[2pt]
    {\itshape\grothFoldertitle}\\[4pt]
    {\small feuillets \grothFirst--\grothLast\ (lot \grothBatch)}\\[2pt]
    {\footnotesize\color{gray!60!black}Transcription établie d'après le fac-similé numérisé
     par l'Université de Montpellier.\\
     Les lectures incertaines sont soulignées, les illisibles notées [\,\ldots\,],
     les ajouts entre crochets.}
  \end{center}
  \vspace{1em}}

\endinput


\folder{101}
\batch{2}
\pages{21}{40}
\dating{[à partir de 1967-1968]}
\watermark{Édition de démonstration}
\foldertitle{Cohomologie locale : notes manuscrites (s.d.).}

\begin{document}

\section{[Modules fondamentaux et dualité locale (suite)]}
\note{titre de l'éditeur, entre crochets ; le feuillet 21 reprend une
démonstration commencée au lot précédent.}

\page{21}
\uncertain{Que} $\Omega$ soit Coh.\ Mac.\ de dim $n$ résulte des formules
précédentes et de la caractérisation des modules de Coh.\ Mac.\ de dim $n$.
Il s'ensuit \uncertain{et par \ill{} d'après le Th.\ 3}
\[
  H^i(\Omega) = 0 \quad \text{pour } i < n ,
\]
\struck{\ill{}} \uncertain{D'après} le Th.\ 1, le cor.\ 1 implique que
$H^n(\Omega)$ est dualisant.
\note{« et par \ill{} d'après le Th. 3 » est écrit en interligne au-dessus
de la fin de la ligne ; « D'après le Th. 1, » est encadré par un trait qui
le rattache à ce qui suit.}

D'après le th.\ de dualité, \ill{}
\[
  H^i(\Omega) \simeq D\bigl(\mathrm{Ext}^{n-i}(\Omega, \Omega)\bigr)
\]
\ill{} \uncertain{qui implique} que \struck{$H^i(\Omega) = 0$ si $i > n$}
\[
  \mathrm{Ext}^{n-i}(\Omega, \Omega) = 0 \quad \text{si } \uncertain{i \neq n},
  \ \text{i.e. } n - i \neq 0 .
\]
\struck{\ill{} \ill{} \ill{} \ill{}} Si $n - i = 0$, on trouve
\[
  D\bigl(\mathrm{Hom}(\Omega, \Omega)\bigr) \simeq H^n(\Omega)
\]
\uncertain{qui est} dualisant, i.e.\ \ill{} $\simeq D(A)$, d'où
$\mathrm{Hom}(\Omega, \Omega) \simeq A$. Ainsi \emph{tout endomorphisme de
$\Omega$ est donné par un \uncertain{scalaire}, déterminé de façon
unique}\,; \uncertain{a fortiori}, \emph{$\Omega$ est \struck{\ill{}} un
module fidèle}.
\note{« a fortiori » est écrit en interligne, avec un crochet qui le place
devant « $\Omega$ est un module fidèle ».}

\emph{$A$, $\Omega$, \uncertain{comme} ci-dessus.}

\textbf{Corollaire 4.} $\Omega_A$ est de dim.\ inj.\ $= n$, \uncertain{en}
\uncertain{particulier} \ill{}.
\[
  \mathrm{Ext}^i(M, \Omega) = 0 \quad \text{si } i > n
\]
\emph{pour tout $M$}.

Il suffit en effet de le vérifier pour $M$ de type fini, et \uncertain{ça}
résulte alors du th.\ de dualité.

\marginal{\textbf{Corollaire 5.} Soit $A$ un anneau local de Coh.\ Mac.\ de
dim $n$, $\Omega$ un module fondamental pour $A$, $M$ un $A$-module de type
fini. \uncertain{Alors} \struck{\ill{}} $\mathrm{codim}\, M \leq k
\Longleftrightarrow \mathrm{Ext}^i_A(M, \Omega) = 0$ pour
$i \mathrel{\uncertain{\geq}} n - k$. \uncertain{Si} $M$ est de dim Krull
$k$, \ill{} \ill{} : $M$ \uncertain{soit} Coh.\ Mac.\ \ill{}
$\mathrm{Ext}^i_A(M, \Omega) = 0$ pour $i \neq k$.}
\note{le corollaire 5 est écrit le long du bord gauche du feuillet, feuille
tournée, sur cinq lignes ; « type » est ajouté au-dessus de la ligne. Le
signe entre $i$ et $n - k$ est un $>$ suivi d'un point, peut-être $\geq$.}

\page{22}
$A$ anneau local de dim $n$, $I$ un module dualisant, $\Omega$ un
\emph{module fondamental pour $A$}, avec un isom donné
\[
  \boxed{\mathfrak{D}_1 : H^n(A) \xrightarrow{\ \sim\ } D(\Omega)}
\]
correspondant à un \emph{accouplement}
\[
  \boxed{H^n(A) \otimes \Omega = H^n(\Omega) \xrightarrow{\ \mathfrak{D}\ } I ,
  \qquad \mathfrak{D} \in D\bigl(H^n(\Omega)\bigr)}
\]
\note{en regard, à gauche, « [\ill{} \ill{} » et à droite
« \struck{\ill{}} ] » : un commentaire entre crochets dont ne se lisent que
les bords.}
et à un isom
\[
  \mathfrak{D}_2 : \hat{\Omega} \xrightarrow{\ \sim\ } D\bigl(H^n(A)\bigr)
\]

On a pour tout $M$ \emph{de type fini}
\[
  \boxed{H^n(M) \xleftarrow[\ \sim\ ]{\mathrm{can}} H^n(A) \otimes M}
  \xrightarrow[\ \sim\ ]{\mathfrak{D}_1 \otimes \mathrm{id}_M}
  D(\Omega) \otimes M \xrightarrow[\ \sim\ ]{\mathrm{can}}
  D\bigl(\mathrm{Hom}(M, \Omega)\bigr)
\]
d'où
\[
  \boxed{\begin{array}{l}
    \mathfrak{D}_{1,M} : H^n(M) \xrightarrow{\ \sim\ }
      D\bigl(\mathrm{Hom}(M, \Omega)\bigr) \\
    \mathfrak{D}_{2,M} : \mathrm{Hom}(M, \Omega)^{\wedge}
      \xrightarrow{\ \sim\ } D\bigl(H^n(M)\bigr)
  \end{array}}
\]
et par là, qui pour $M \simeq A$ se réduisent à $\mathfrak{D}_1$,
$\mathfrak{D}_2$. Pour $M = \Omega$, on trouve
\[
  \begin{array}{l}
    \mathfrak{D}_{1,\Omega} : H^n(\Omega) \xrightarrow{\ \sim\ }
      D\bigl(\mathrm{Hom}(\Omega, \Omega)\bigr) \\
    \mathfrak{D}_{2,\Omega} : \mathrm{Hom}(\Omega, \Omega)^{\wedge}
      \xrightarrow{\ \sim\ } D\bigl(H^n(\Omega)\bigr)
  \end{array}
\]
Enfin, on constate que l'élément identique de
$\mathrm{Hom}(\Omega, \Omega)^{\wedge}$ correspond à l'élément
$\mathfrak{D}$ de $D(H^n(\Omega))$. Par suite, le transposé de
$\mathfrak{D}$ est l'hom
\[
  {}^t\mathfrak{D} : \hat{A} \longrightarrow \mathrm{Hom}(\Omega, \Omega)^{\wedge}
\]
qui transforme $1_{\hat{A}}$ en $\mathrm{id}_{\hat{\Omega}}$, \ill{}
$\lambda$ \ill{} l'homothétie $\lambda$.

\textbf{Proposition 9.} L'annulateur de $\Omega$ est l'intersection
$\mathfrak{N}$ des idéaux primaires associés aux idéaux premiers de
\uncertain{corang} $n$ dans $A$, et $\Omega$ est aussi un module
fondamental pour $A' = A/\mathfrak{N}$. Les idéaux premiers associés à
$\Omega$ sont les idéaux pr.\ de \uncertain{corang} $n$ de $A$.
\note{le trait vertical qui encadre l'énoncé à gauche porte en marge :}
\marginal{$\mathfrak{p}$ \ill{} \ill{} ;
$\ell_{\mathfrak{p}}(\Omega) = \ell_{\mathfrak{p}}(A)$ \ill{}
$\mathfrak{p}$}
\note{la lettre $\mathfrak{N}$ transcrit une majuscule gothique de lui,
lue ainsi dans tout le lot ; « corang » est une lecture incertaine d'un mot
bref devant « $n$ », tenue la même aux pages suivantes.}

L'annulateur de $\Omega$ est celui de $D(\Omega) = H^n(A)$, \uncertain{donc}
\struck{\ill{}} la première assertion suit du th.\ 3, la

\page{23}
dernière suit du cor.\ 1 dudit. On a
\[
  0 \longrightarrow \mathfrak{N} \longrightarrow A \longrightarrow A'
  \longrightarrow 0
\]
d'où, comme $\mathfrak{N}$ est de dim $< n$ donc $H^n(\mathfrak{N}) = 0$,
\[
  H^n(A) \simeq H^n(A')
\]
donc $\Omega$ est aussi fond.\ pour $A'$.

[N.B. $\mathfrak{N}$ est le plus grand idéal \uncertain{de} $A$ \ill{}
\uncertain{dimension} $< n$. \ill{} que $\Omega$ aussi \ill{} \ill{}
\ill{}]

\textbf{Corollaire 1.} L'image de $H^n(\Omega) \xrightarrow{\mathfrak{D}} I$
est l'annulateur de $\mathfrak{N}$,
\uncertain{i.e.\ $\mathrm{Coker}\,\mathfrak{D} \simeq D(\mathfrak{N})$}
\struck{Donc \ill{} $\mathfrak{D}$ est surjectif \ill{} \ill{}}
\struck{si \ill{} \ill{} \ill{} les idéaux}
\note{« i.e. Coker $\mathfrak{D} \simeq D(\mathfrak{N})$ » est en
interligne, souligné ; les deux lignes biffées étaient une première
rédaction de la liste qui suit.}
\begin{enumerate}
  \item[(i)] $\mathfrak{D}$ est surjectif
  \item[(ii)] $\mathfrak{N} = 0$, i.e.\ les idéaux premiers associés à $0$
    dans $A$ sont \struck{\ill{}} de \uncertain{corang} $n$
  \item[(iii)] $\Omega$ est un $A$-module fidèle.
\end{enumerate}

En effet, la première assertion résulte \uncertain{du fait} que
$\mathfrak{N}$ est le noyau de $A \to \mathrm{Hom}(\Omega, \Omega)$, qui
est \uncertain{transposé} de $\mathfrak{D}$. D'où les équivalences.

\textbf{Corollaire 2.}
$\mathrm{Ker}\,\mathfrak{D} \simeq D\bigl(\mathrm{Hom}(\Omega, \Omega) /
A \cdot \mathrm{id}_{\Omega}\bigr)$. \emph{Les conditions suivantes sont
équivalentes} :
\begin{enumerate}
  \item[(i)] $\mathfrak{D}$ est injectif \qquad (i bis) $H^n(\Omega)$ est
    \uncertain{comonogène}
  \item[(ii)] Tout endom.\ de $\Omega$ est une homothétie
  \item[(ii bis)] $\mathrm{Hom}(\Omega, \Omega)$ est monogène.
\end{enumerate}

La première formule \ill{} \ill{} dualité.
$(\mathrm{i}) \Leftrightarrow (\mathrm{ii\,bis})$,
\struck{\ill{}} $(\mathrm{ii\,bis}) \Leftrightarrow (\mathrm{i\,bis})$ et

\page{24}
$(\mathrm{ii}) \Longrightarrow (\mathrm{ii\,bis})$ trivialement. Enfin
$(\uncertain{\mathrm{i\,bis}, \mathrm{ii\,bis}}) \Longrightarrow
(\mathrm{ii})$, car si $M$ est un module de type fini sur $A$ tel que
$\mathrm{Hom}(M, M)$ \uncertain{soit} monogène, il résulte de Nakayama que
$\mathrm{id}_M$ en est un générateur.

\textbf{Cor.\ 3.} Conditions équivalentes
\begin{enumerate}
  \item[(i)] $\mathfrak{D}$ est bijectif
  \item[(ii)] $H^n(\Omega)$ est dualisant
  \item[(iii)] $A \to \mathrm{Hom}(\Omega, \Omega)$ est un isom
  \item[(iv)] $\mathrm{Hom}(\Omega, \Omega) \simeq A$
\end{enumerate}
\struck{\ill{}} Alors \struck{\ill{}} les idéaux premiers associés à $A$
sont tous de \struck{dim} \uncertain{corang} $n$.

\textbf{Cor.\ 4.} Les cond.\ du cor.\ 3 sont vérifiées dans les conditions
suivantes : (a) $A$ est Cohen-Macaulay \qquad (b) $A$ est
\uncertain{normal}.

(a) Résulte du th.\ \uncertain{4}, cor.\ 3.

Dans le cas (b), \uncertain{on} \ill{} la condition du cor.\ 1 vérifiée
($A$ \uncertain{étant} intègre) il suffit \uncertain{de prouver} que
$\mathrm{Hom}(\Omega, \Omega) \supset A$ est \ill{} \ill{}. Or $\Omega$
est sans torsion \ill{} \ill{} \ill{} idéal premier,
[\struck{\ill{} \ill{}} \uncertain{p.\ rapport} \ill{} $A$ intègre] donc
$\mathrm{Hom}(\Omega, \Omega) \subset K$ corps
\marginal{à \uncertain{expliciter} \ill{}}
\note{« [\ill{} p. rapport … intègre] » est une incise entre crochets,
dont le premier mot est biffé et dont la fin est écrite en interligne ; la
note marginale, oblique, lui est reliée par une accolade.}

\page{25}
des fractions de $A$. \struck{Comme} c'est un module de type fini sur $A$,
il est dans la clôture intégrale de $A$, O.K.

\textbf{Question.} Les conditions du cor.\ 2 sont-elles \emph{tjrs}
vérifiées\,? Cela est \uncertain{vraisemblable}. \ill{} \ill{} les idéaux
pr.\ associés à $0$ dans $A$ sont de \struck{\ill{}} \uncertain{corang}
$n$ : est-il vrai \ill{} que $H^n(\Omega)$ est dualisant\,?
\emph{Non} \ill{} \ill{} \struck{\ill{} \ill{}} \ill{} \ill{} dans le cas
\struck{\ill{} \ill{} \ill{} \ill{} \ill{}} \uncertain{loc.} 2
\uncertain{intègre} $A$ [de dim 2], tel que $A_{\mathfrak{p}}$ soit
\uncertain{premier unibranche} pour $\mathfrak{p}$ \struck{\ill{}}
[$A$ \ill{} \ill{}]. En effet, \struck{\ill{}} on vérifie \ill{} \ill{} que
\[
  \mathrm{Hom}_A(\Omega, \Omega) = \bigcap_{\substack{\mathfrak{p}\ \text{premier}\\
  \text{de corang 1}}} A_{\mathfrak{p}} =
  \struck{\text{normalisé}}\ \text{clôture intégrale de } A .
\]
[$\Omega$ \ill{} \ill{} \ill{} \ill{} conditions \ill{} \ill{}.]
\note{sous l'intersection, « corang » est une lecture incertaine, comme
partout dans le lot. La question est encadrée à gauche d'un trait en zigzag ; plusieurs
mots en interligne au-dessus de « dans le cas » ne se lisent pas ;
« intègre » est en interligne avec une accolade, « premier unibranche »
en interligne au-dessus du mot biffé. « clôture intégrale » est écrit
au-dessus de « normalisé » biffé.}

\textbf{Proposition 10.} Soit $x$ un élément de $A$ qui ne soit pas diviseur
de $0$, et tel que $x H^{n-1}(A) = H^{n-1}(A)$. Alors
$\Omega_A / x\Omega_A$ est un module fondamental pour $A/xA$.
\note{« Alors » est d'une encre plus pâle ; dans « $A/xA$ », deux lettres
repassées d'une encre plus noire se superposent à une première écriture.}

\page{26}
En effet, on a la suite exacte
\[
  0 \longrightarrow A \xrightarrow{\ x\ } A \longrightarrow A/xA
  \longrightarrow 0
\]
d'où la suite
\[
  H^{n-1}(A) \xrightarrow{\ x\ } H^{n-1}(A) \longrightarrow H^{n-1}(A/xA)
  \longrightarrow H^n(A) \xrightarrow{\ x\ } H^n(A)
\]
et compte tenu de l'hypothèse
\[
  H^{n-1}(A/xA) \simeq \struck{\mathrm{Ker}\ \text{de}\ x\ \text{dans}}\
  \text{noyau de } x \text{ dans } H^n(A) .
\]
Comme $H^n(A) \simeq D(\Omega)$, on trouve bien
\[
  H^{n-1}(A/xA) \simeq D(\Omega/x\Omega) , \quad \text{cqfd.}
\]

\textbf{Corollaire.} Supposons $A$ Coh.\ Mac.\ et soit $x_1, \ldots, x_p$
une partie d'un syst.\ de paramètres de $A$. Alors
$\Omega_A / (x_1, \ldots, x_p)\Omega_A$ est fondamental pour
$A/(x_1, \ldots, x_p)$.

Supposons maintenant que $A$ soit un quotient d'un anneau local
\struck{\ill{}} Coh.\ Mac.\ $A'$ \struck{\ill{}} :
\[
  A \simeq A'/\mathfrak{N}
\]
\uncertain{Supposons} $A'$ \uncertain{muni} d'un $\Omega_{A'}$. Si $M$ est
un $A$-module, c'est un $A'$-module, et $D(M)$ \uncertain{peut}
s'interpréter \uncertain{ainsi}. [\uncertain{divisons} un module dualisant
$I'$ ($= H^{n'}(\Omega_{A'})$, pour $A'$) \ill{} \ill{} que
$\Omega_{A'}$]. \ill{} \ill{} \ill{}
\note{« Coh. Mac. » est écrit au-dessus d'un mot biffé, peut-être
« régulier » ; l'exposant de $H^{n'}$ est repassé. La page s'arrête sur
une phrase coupée.}

\page{27}
i.e.\ \ill{} la \uncertain{forme} d'\uncertain{application} \ill{} la
\uncertain{th.}\ de dualité à $A$, l'\uncertain{application} à $A'$,
\ill{} \ill{}
\[
  (1) \qquad H^i(M) \simeq
  D\bigl(\mathrm{Ext}^{\ill{}-i}_{A'}(M, \Omega_{A'})\bigr)
  \qquad \text{i.e.}
\]
\struck{(1 bis) \ill{} $\mathrm{Ext}^{n'-i}_{A'}$ \ill{}}
\note{l'exposant de $\mathrm{Ext}$ est noirci devant « $-i$ » ; la ligne
(1 bis) est biffée d'un trait ondulé.}

\uncertain{On} \ill{} \uncertain{dim} \ill{} \uncertain{particulier}
\[
  (2) \qquad H^n(A) \simeq
  D\bigl(\mathrm{Ext}^{n'-n}_{A'}(M, \Omega_{A'})\bigr)
\]
\ill{} \uncertain{qui} \uncertain{permet} \ill{} \uncertain{de}
\uncertain{poser} \struck{\ill{}} \uncertain{par} \uncertain{définition}
\[
  (2\ \text{bis}) \qquad \Omega_A \simeq
  \mathrm{Ext}^{n'-n}_{A'}(M, \Omega_{A'}) .
\]
\note{dans (2) et (2 bis), et plus loin dans (13), l'argument de
$\mathrm{Ext}$ est écrit $M$ là où l'on attend $A$ ; on garde la lettre
de la page.}

\struck{Plus généralement, \ill{}, \ill{}}
\struck{(3) $\Omega$ \ill{}}
\note{ces deux lignes sont encadrées et biffées.}

Soit $\mathfrak{p}$ un idéal premier de $A$, \ill{}
$\mathfrak{p} \simeq A'/\mathfrak{p}'$,
\uncertain{où} $\mathfrak{p}' \supset \mathfrak{N}$. \uncertain{On a}
\uncertain{alors}
\[
  (\Omega_A)_{\mathfrak{p}} \simeq
  \mathrm{Ext}^{n'-n}_{A'_{\mathfrak{p}'}}\bigl(M_{\mathfrak{p}^*},
  (\Omega_{A'})_{\mathfrak{p}'}\bigr) .
\]
\note{« $\mathfrak{p} \simeq A'/\mathfrak{p}'$ » tel qu'écrit ; l'indice
« $\mathfrak{p}^*$ » de $M$ aussi.}

Supposons \uncertain{qu'on} \uncertain{sache} que
$(\Omega_{A'})_{\mathfrak{p}'}$ est un module fond.\ pour l'anneau
[\emph{Coh.\ Mac.}] $A'_{\mathfrak{p}'}$ : ce sera le cas en particulier si
$A'$ est \emph{régulier}, car alors $\Omega_{A'} \simeq A'$, et
$A'_{\mathfrak{p}'}$ sera régulier (Serre). \ill{} \ill{} de
$A_{\mathfrak{p}}$, \ill{} \ill{} un module fond. Alors
$(\Omega_A)_{\mathfrak{p}}$ \ill{} un module dualisant pour
$A_{\mathfrak{p}}$, \ill{} (2 bis) \uncertain{appliqué} à
$A_{\mathfrak{p}} \uncertain{\simeq}
A'_{\mathfrak{p}'}/\mathfrak{N}A'_{\mathfrak{p}'}$.
\note{« dualisant » est écrit là où l'on attend « fondamental » ; devant
$A_{\mathfrak{p}}$ une première lettre est surchargée.}

\page{28}
D'où :

\textbf{Prop.\ 11.} Soit $A$ un quotient d'un anneau local régulier,
\uncertain{alors} $\Omega$ est un module fond.\ pour $A$. Si
$\mathfrak{p} \subset A$ est premier dans $A$, $\Omega_{\mathfrak{p}}$ est
fond.\ pour $A_{\mathfrak{p}}$.

[N.B. Une \uncertain{démonstration} \uncertain{directe} \ill{}
\uncertain{question} de définition, \uncertain{sans}
\uncertain{interprétation} \ill{}]. De plus,
\ill{} \uncertain{traduire} \ill{} (1) \ill{}

\textbf{Lemme.} Pour que \uncertain{$\mathrm{codim}$} $M \leq k$, il f.\
et s.\ que
\[
  \mathrm{Ext}^{n'-i}_{A'}(M, \Omega_{A'}) = 0 \quad \text{pour}
  \ 0 \leq i < k .
\]
[Fait bien \uncertain{connu} \ill{}, \uncertain{dim} : \ill{}].
\struck{\ill{}}

\textbf{Corollaire.} $A$ Coh.-Mac.\ $\Longleftrightarrow$
$\mathrm{Ext}^{n'-n+i}_{A'}(M, \Omega_{A'}) = 0$ pour $i > 0$.

\uncertain{Cor.} \ill{}, \ill{} : \uncertain{introduisons} \ill{}
\[
  (13) \qquad \Omega_A^{(i)} =
  \mathrm{Ext}^{n'-n+i}_{A'}(M, \Omega_{A'}) :
\]
\ill{} (\ill{} isom \uncertain{non canon.}\ près) \ill{}
\[
  (13\ \text{bis}) \qquad D(\Omega_A^{(i)}) \simeq H^{\uncertain{n-i}}(A)
\]
[N.B. \uncertain{Analogie} \ill{} \ill{} \ill{} \ill{} \ill{} des
$\Omega_A^{(i)}$].
\note{les exposants de $\mathrm{Ext}$ sont écrits « $n'-i$ »,
« $n'-n+i$ » avec un premier signe repassé à l'encre noire ; celui de
(13 bis) est noirci.}

\uncertain{D'où} :
\[
  \Omega_A^{(0)} = \Omega_A
\]
\uncertain{mod.}\ fond., et
\[
  \Omega_A^{(i)} = 0 \ \text{pour} \ i \geq 1 \Longleftrightarrow A \
  \text{est Coh.\ Mac.}
\]

\page{29}
On peut \uncertain{expliquer}, en le \uncertain{généralisant}, le th.\ de
dualité locale par une suite spectrale.

\textbf{Lemme 6.} $A' \to A$ un hom d'anneaux, $N$ un $A$-module, $M$ un
$A'$-module. Il y a une suite spectrale \uncertain{cohomologique}
\[
  \mathrm{Ext}^{\bullet}_{A'}(N, M) \Longleftarrow
  \mathrm{Ext}^p_A\bigl(N, \mathrm{Ext}^q_{A'}(A, M)\bigr)
\]
\note{dans cette formule et les suivantes, plusieurs $A$ sont repassés à
l'encre noire sur une première lettre, qui semble un $B$ ; un $B$ non
corrigé subsiste dans $\mathrm{Hom}_{A'}(B, M)$ ci-dessous.}

\ill{} \ill{} \ill{} \ill{} foncteur d'\ill{} :
\[
  \mathrm{Hom}_{A'}(N, M) = \mathrm{Hom}_A\bigl(N, \mathrm{Hom}_{A'}(B, M)\bigr)
\]
\[
  \mathcal{C}^{A'} \xrightarrow{\ h^{A'}_A\ } \mathcal{C}^{A}
  \xrightarrow{\ h^A_N\ } \mathcal{G}^{\uncertain{A}} ,
  \qquad h^A_N \circ h^{A'}_A = h^{A'}_N
\]
\note{la composée est tracée comme un arc sous les deux flèches et
étiquetée $h^{A'}_N$ ; les exposants des trois lettres rondes sont
repassés et ne se lisent pas avec sûreté.}

$h^{A'}_A$ transforme injectifs en injectifs.

- - - - -

\textbf{Prop.\ 12.} \marginal{Soit $A$ un anneau quotient d'un Coh.\ Mac.\
$A'$ muni d'un module fond.\ $\Omega_{A'}$.} Il y a une suite spectrale
canonique de type cohomologique
\[
  \boxed{\bigl(E^{i}(M)\bigr) = \bigl(H^{\ill{}i}(M)\bigr) \Longleftarrow
  E_2^{pq}(M) = D\bigl(\mathrm{Ext}^p_A(M, \Omega_A^{(q)})\bigr)}
\]
\note{l'hypothèse sur $A$ est écrite au-dessus de l'énoncé, dans une
boucle qui la rattache à « Prop. 12 ». Un signe noirci précède l'indice
$i$ de $H$, un autre l'exposant de $E$ ; on ne les lit pas.}

Cela \uncertain{résulte aussitôt} du lemme 6 avec \uncertain{$N$ remplacé
par} $M$, et $M$ remplacé par $\Omega_{A'}$.

\marginal{\textbf{Cor.\ 1.} Soit $x \in A$ non diviseur de $0$ dans $A$,
ni dans $M$. Alors \ill{} $\mathrm{Ext}^p_A(N, M) \simeq
\mathrm{Ext}^{p-1}_{A/x}(N, M/xM)$ pour tout $N$, \ill{} \ill{}.
\textbf{Cor.\ 2.} \ill{} \ill{} $M$-\ill{}}
\marginal{\emph{Modules fondamentaux} \emph{pour un} \ill{} \ill{}
\uncertain{anneau} \ill{} \ill{}}
\note{le premier bloc est écrit le long du bord gauche, feuille tournée ;
le second en oblique au pied de la marge, souligné mot à mot, comme une
annonce de ce qui suit.}

\section{6. Caractérisation des modules fondamentaux d'anneaux Coh.\ Mac. Anneaux de Gorenstein.}
\note{titre de sa main, en tête du feuillet 30, souligné.}

\page{30}
\note{toute la page, jusqu'à « c') » inclus, est barrée d'une grande croix ;
la « Proposition 13 » du pied de page, qui en reprend l'énoncé, ne l'est
pas.}

\struck{\textbf{Théorème 5.} Soient $A$ un anneau local noeth., $M$ un}
\struck{$A$ \ill{}}

\textbf{Prop.\ 13.} Soient $A$ un anneau local noeth., \uncertain{de corps
rés.\ $k$,} $M$ un $A$-module, $n$ un entier $> 0$, \ill{}
\[
  \mathrm{Ext}^{n-1}(k, M) = 0 .
\]
Alors les conditions \uncertain{suivantes} sont équivalentes
\begin{enumerate}
  \item[a)] $\mathrm{Ext}^n(k, M) = k$ et $\mathrm{Ext}^{n+1}(k, M) = 0$
  \item[b)] $N \mapsto \mathrm{Ext}^n(N, M)$ est dualisant
  \item[c)] $H^n(M)$ est un module dualisant.
\end{enumerate}
[\uncertain{Si} \ill{} $n = \dim A$, les idéaux pr.\ associés à $A$ sont
\uncertain{de dim $n$}]
\struck{et $\mathrm{lg}_{\mathfrak{p}}(M) = \mathrm{lg}_{\mathfrak{p}}(A)$
pour \ill{} $\mathfrak{p}$} \uncertain{$n = \dim A$.}

\textbf{Cor.\ 1.} Soit $\mathfrak{p}$ un idéal pr.\ minimal de $A$, les
conditions \uncertain{précédentes} \struck{équivalent aux conditions :
$n = \dim A$,} \struck{\ill{} \ill{}}
\[
  \struck{0 < \mathrm{lg}_{\mathfrak{p}} M \leq \mathrm{lg}_{\mathfrak{p}} A}
  \quad \text{et} \ \ill{} \qquad
  \boxed{0 < \mathrm{lg}_{\mathfrak{p}} M \leq \mathrm{lg}_{\mathfrak{p}} A}
\]
\ill{} \ill{} \uncertain{conditions} \uncertain{équivalentes}
\uncertain{suivantes}
\note{l'inégalité, d'abord écrite puis biffée dans la ligne, est récrite
dans une boucle qu'une flèche relie à « $n = \dim A$ » de la ligne
au-dessus ; un mot est écrit en interligne au-dessus de la ligne biffée.}

\struck{$\mathrm{Ext}^{n-1}(k, M) = 0$ \ill{} \ill{}}

à \ill{} \ill{} \ill{}
\begin{enumerate}
  \item[a')] $\mathrm{Ext}^{n+1}(k, M) = 0$
  \item[b')] $N \mapsto \mathrm{Ext}^n(N, M)$ est exact
  \item[c')] $H^n(M)$ est injectif
\end{enumerate}

\textbf{Proposition 13} \struck{Lemme 7}. Soit $A$ \ill{} local,
\uncertain{de dim $n$}, $M$ un $A$-module, \struck{\ill{}}
\struck{[i.e.\ $\mathrm{Ext}^n(N, M) \simeq \mathrm{Hom}(N, H^n(M))$]}
\struck{Alors les conditions \ill{}} \ill{}
\note{« Proposition 13 » est écrit au-dessus d'un « Lemme 7 » biffé ;
« de dim $n$ » est en interligne, avec un crochet qui le place après
« local ».}

\uncertain{Supposons}
\[
  \boxed{\mathrm{Ext}^{n-1}_A(k, M) = 0}
\]
\struck{\ill{}} \ill{} \uncertain{les} \uncertain{suivantes sont}
\uncertain{équivalentes}
\begin{enumerate}
  \item[a')] $\mathrm{Ext}^{n+1}(k, M) = 0$
  \item[b')] $N \mapsto \mathrm{Ext}^n(N, M)$ est exact ($N$ de type fini)
  \item[c')] $H^n(M)$ est injectif
\end{enumerate}

\page{31}
\uncertain{Démonstration} :
$\mathrm{a} \Rightarrow \mathrm{b} \Rightarrow \mathrm{c}$ trivial,
$\mathrm{c} \Rightarrow \mathrm{a}$ car $N \mapsto \mathrm{Ext}^n(N, M)$
\uncertain{est exact}, donc
$N \mapsto \mathrm{Ext}^{n+1}(N, M)$ \uncertain{est exact} : \ill{},
donc $\mathrm{Ext}^{n+1}(k, M) \to H^{n+1}(M)$ est injectif, et
$H^{n+1}(M) = 0$.

\textbf{Cor.\ 1.} \ill{}, \ill{} \ill{} \ill{} \ill{} conditions,
[les idéaux \ill{}
\struck{\ill{} \ill{} \ill{} \ill{} \ill{}} $M$,] \ill{} \ill{} de dim
$n$, [et \ill{} \ill{} \ill{} \ill{}]
\[
  H^n(M) \simeq D(A)^m \ \text{pour un entier} \ m
\]
\struck{où} \uncertain{dans} $\uncertain{\sqrt{A}}$ \ill{} \ill{}, \ill{}
\uncertain{tout} $\mathfrak{p}$
\[
  \mathrm{lg}_{\mathfrak{p}}(M) = m\, \mathrm{lg}_{\mathfrak{p}}(A)
  \qquad
  \begin{array}{l}
    \mathrm{Ext}^n(k, M) \simeq k^m \\
    \mathrm{Ext}^i(k, M) = 0 \ \text{si} \ i > n
  \end{array}
\]
\note{l'énoncé est très repris ; plusieurs incises entre crochets, dont
une ligne biffée, sont écrites en interligne à droite.}

\struck{\textbf{Cor.} \ill{} \ill{} \ill{}
$[0 < \mathrm{lg}_{\mathfrak{p}}(M) \leq \mathrm{lg}_{\mathfrak{p}}(A)]$
les conditions \ill{} \ill{} \ill{}}
\struck{\ill{} $\mathrm{Ext}^n(M, M)$ \ill{} $H^n(M)$ est dualisant}
\note{deux lignes barrées de plusieurs traits.}

\textbf{Remarque.} [\ill{} \ill{} : 2 \ill{} \ill{} \ill{}]
($\Omega$ \uncertain{satisfait} aux conditions \ill{})

\ill{} \struck{\ill{}} $A$ \uncertain{est} Coh.-Mac.\ et
\struck{\ill{} \ill{}}
\[
  M \simeq \Omega^m , \ \text{où} \ \Omega \ \text{est un module
  fondamental}
\]
\uncertain{pour} $A$, [et \uncertain{où} l'hyp.\
$\mathrm{Ext}^{n-1}(k, M) = 0$ est \uncertain{=} \uncertain{redondante}
\uncertain{dans} b')] \ill{}, \uncertain{précisément} \uncertain{pour}
\ill{} \ill{}, il \ill{} \ill{} \ill{} \ill{} \uncertain{pour} $n > 0$,
\ill{} : \ill{} \ill{} \ill{} \ill{} $M$ \ill{} \ill{}. \uncertain{Voir}
\ill{}, \ill{} \ill{} \ill{} \ill{} \ill{}
\note{ce passage est serré et rapide ; ses mots de liaison ne se lisent
pas. Une incise entre crochets est écrite au-dessus de « pour $A$ ».}

\marginal{(\uncertain{Supposons} $m = 1$, \struck{\ill{}} $A$ \ill{}
\uncertain{quotient} \uncertain{régulier}.)}
\struck{\uncertain{Sous} les conditions \uncertain{précédentes},}
\note{la parenthèse est entourée d'une boucle.}

\textbf{Proposition 14.} \uncertain{Supposons} \ill{} \ill{} \ill{}
\ill{} conditions a) b) c) \struck{\ill{} \ill{} \ill{}} \ill{} :
$M$ Coh.\ Mac.\ \ill{} $A$ \uncertain{est} Coh.\ Mac.\ et
$M \simeq \Omega_A^m$, où $\Omega_A$ est un module fondamental pour $A$.

C'est trivial si $n = 0$, \uncertain{procédons} \uncertain{par}
\uncertain{réc.}\ sur $n$. \uncertain{Voir} les \ill{} \ill{} - - - -

\marginal{\textbf{Cor 2.} Les conditions précédentes, jointes à :
$0 < \mathrm{lg}_{\mathfrak{p}}(M) \leq \mathrm{lg}_{\mathfrak{p}}(A)$,
\uncertain{sont} équivalentes aux conditions suivantes
\uncertain{compte tenu de} $\mathrm{Ext}^{n-1}(k, M) = 0$ :
a) $\mathrm{Ext}^n(k, M) = k$, $\mathrm{Ext}^{n+1}(k, M) = 0$ ;
b) $N \mapsto \mathrm{Ext}^n(N, M)$ \uncertain{est} \uncertain{dualisant} ;
c) $H^n(M)$ \uncertain{est} \uncertain{dualisant}.
[\uncertain{équivalentes}]}
\marginal{\textbf{Cor 3.} Les conditions du cor.\ 2 sont
\uncertain{entraînées} par l'\uncertain{une} des cond.\ \uncertain{suivantes} :
d) $\mathrm{Ext}^n(k, M) = k$, $M$ fidèle ;
d bis) $\mathrm{Ext}^n(A/\mathfrak{N}, M) \to$
\struck{\ill{}} \uncertain{comonogène}, $M$ fidèle.
\uncertain{Pour} \struck{\ill{}} \struck{\ill{}} \ill{} aux idéaux
\uncertain{premiers} \ill{} $M$ \ill{} \ill{} \ill{} \ill{}.}
\note{les corollaires 2 et 3 sont écrits le long du bord gauche, feuille
tournée, chacun encadré d'un trait ; le « $\mathfrak{N}$ » de d bis) est
mal formé.}

\page{32}
\struck{Proposition 15}

\struck{\ill{} \ill{}}

\textbf{Cor.\ 1.} \struck{\ill{} \ill{} \ill{}} \uncertain{Supposons}
\ill{} $M$ Coh.\ Mac.\ sur $A$ de dim $n$
\struck{\ill{} \ill{} \ill{}, \ill{} \ill{}}
(i.e.\ $A$ est Coh.\ Mac.\ et $M \simeq \Omega_A$ :)
Alors les conditions \struck{\uncertain{précédentes}} \ill{} \ill{}
\uncertain{équivalentes}

a) b) c) du cor.\ 2 à prop.\ 13 ; d / d bis) du cor.\ 3
\begin{enumerate}
  \item[\struck{d)}] \struck{$\mathrm{Ext}^n(k, M) = k$, $M$ \ill{}
    fidèle}
  \item[\struck{e)}] \struck{$\mathrm{Ext}^n(A/\mathfrak{N}, M)$
    \ill{} \uncertain{comonogène}, $M$ \ill{} fidèle}
  \item[\struck{f)}] \struck{\ill{} \ill{} \ill{} \ill{}}
\end{enumerate}
\note{d), e) et f) sont encadrés et biffés ligne à ligne, et traversés de
longs traits obliques ; une boucle ramène « soit $(x_1 \ldots x_n)$ un
syst. de paramètres de $M$ », écrit à droite, dans ce qui suit.}

\struck{\textbf{Cor.} \ill{} \ill{} \ill{}}

$x_1 \ldots x_n$ un syst.\ \uncertain{de paramètres} de $M$. \uncertain{Alors}
\ill{} les conditions \struck{précédentes équivalent} à \uncertain{la}
\begin{enumerate}
  \item[f)] $M/x_1M + \cdots + x_nM$ \struck{est} est
    \uncertain{comonogène}.
\end{enumerate}

\textbf{Remarque.} Il est possible que dans le cor.\ 1, on
\uncertain{puisse} \uncertain{remplacer} \struck{Ext} la cond.\ : $M$
Coh.\ Mac.\ par la condition :
$\mathrm{Ext}^{n-1}(k, M) = 0$. Il \uncertain{semble} $=$ \ill{} \ill{}
\uncertain{la} condition \uncertain{soit} \uncertain{nécessaire} dans
d bis) \struck{\ill{} \ill{} \ill{} e)} \uncertain{des}
[\uncertain{et} \uncertain{implique} \uncertain{les} b), c) \uncertain{et}
d)] \ill{} \ill{} \ill{} \ill{} f.

La proposition suivante est de nature plus \uncertain{élémentaire}.

\textbf{Proposition 15.} Soit $A$ un anneau local noeth. Les conditions
suivantes sont équivalentes.
\begin{enumerate}
  \item[(a)] $A$ est un module fondamental sur $A$
  \item[(b)] $H^n(A)$ est dualisant
  \item[(c)] $H^n(A)$ est injectif
  \item[(d)] $H^n(A)$ est comonogène et $A$ est
    \uncertain{équidimensionnel}.
\end{enumerate}
\note{les lettres (a), (b), (c), (d) sont repassées sur de premières
lettres.}

\page{33}
(i) \uncertain{signifie} que $D(A) \simeq H^n(A)$, i.e.\ (ii).
(ii) $\Rightarrow$ (iii), d'autre part (iii) $\Rightarrow$
$H^n(A) \simeq D(A)^m$ i.e.\ $A^m$ est \struck{dualisant}
fondamental. En \uncertain{vertu} de prop.\ 9, \ill{}
$\mathrm{lg}_{\mathfrak{p}}(A^m) = \mathrm{lg}_{\mathfrak{p}}(A)$
d'où $m = 1$, d'où (iii) $\Rightarrow$ (i).
\struck{Il \ill{} \ill{} \ill{}} \uncertain{Enfin}
(i) $\Rightarrow$ (iv) par prop.\ 9, réciproquement : $H^n(A)$
\uncertain{comon.}\ \ill{} $\subset D(A)$, \ill{} \ill{}
\ill{} de $H^n(A)$ est \ill{} l'\ill{} \ill{} l'hyp.\ \ill{}
$H^n(A) \simeq D(A)$.

\struck{\textbf{Définition 2.} Un tel $A$ est dit anneau de Gorenstein}
\marginal{\struck{s'il est de plus Cohen-Macaulay}}
\note{la définition est encadrée et biffée ; « s'il est de plus
Cohen-Macaulay » est écrit dessous dans une boucle, biffée avec elle.}

\struck{\textbf{Th.\ 5.} Un anneau de Gorenstein est Coh.\ Mac.}

\struck{En effet, il est \uncertain{équivalent} \ill{}, \ill{} si
\ill{} \ill{}. \uncertain{Réc.}\ sur $n = \dim A$, \uncertain{trivial}
si $n = 0$. Si $n > 0$, \ill{} $A$, \ill{} \ill{}, il y a un
$x \in A$ \uncertain{non diviseur} \ill{} \ill{}. Alors
\ill{} \ill{} \ill{}}
\[
  \struck{0 \to A \xrightarrow{\ x\ } A \to A/xA \to 0}
\]
\[
  \struck{\ill{} \to H^{n-1}(\ill{}) \to H^n(\ill{})}
  \qquad
  \struck{\mathrm{Ext}^{\bullet}(N, A/x) \Longleftarrow
  \mathrm{Ext}^{\bullet}_{A/x}\bigl(N, \mathrm{Ext}^{\bullet}_A(A/x, A)\bigr)}
\]
\struck{d'après le cor.\ \ill{} lemme 6,} \ill{} \ill{}
\note{la moitié inférieure de la page, depuis « Définition 2 », est
biffée de longs traits obliques et de traits horizontaux ligne à ligne,
et reprise de boucles qui en déplacent des morceaux ; on n'en donne que ce
qui se lit.}

\textbf{Prop 16.} (\uncertain{$+$} $A$ est Cohen-Macaulay), les
conditions \uncertain{de la prop.\ 15} \uncertain{sont}
\uncertain{équivalentes} aussi aux suivantes
\begin{enumerate}
  \item[a')] \struck{\ill{}} le foncteur $N \mapsto \mathrm{Ext}^n_A(N, A)$
    \struck{est} \uncertain{dualisant} \struck{\ill{}}
    [\ill{} \ill{} \uncertain{et} exact]
  \item[c')] \ill{}
  \item[\ill{}] $\mathrm{Ext}^n(k, A) \simeq k$ ;
    $\mathrm{Ext}^{n+1}(k, A) = 0$
\end{enumerate}
\note{la fin de la page est surchargée : les lettres a'), c') et une
troisième, noircie, sont reprises, et deux flèches et une longue boucle
déplacent les conditions l'une par rapport à l'autre ; l'ordre donné ici
est celui de la page.}

\page{34}
\begin{enumerate}
  \item[b')] $\mathrm{Ext}^n(N, A)$ dualisant en $N$
  \item[c')] $\mathrm{Ext}^n(N, A)$ exact en $N$
  \item[c'')] $\mathrm{Ext}^{n+1}(\uncertain{k}, A) = 0$, et
    $\mathrm{Ext}^{n-1}(k, A) = 0$
  \item[d)] $\mathrm{Ext}^n(k, A) = k$ et \struck{$\mathrm{Ext}$ \ill{}}
    \uncertain{$A$} est Coh.\ Mac.
\end{enumerate}
\marginal{ici il est facile de voir que ça implique \uncertain{déjà}
Coh.\ Mac.}
\note{la remarque est écrite à droite de b') et c'), qu'une accolade
réunit.}

\struck{\ill{} $A/(x_1, \ldots, x_n)$}

\struck{Immédiat sans rien \ill{} \ill{} \ill{} \ill{} \ill{}
(\ill{} $A$ \ill{} \uncertain{équidimensionnel}) \ill{}, \ill{}}
\note{deux lignes et demie biffées, reprises d'une boucle.}

- - - - - -

\textbf{Corollaire 1.} Ce qui équivaut aussi à :
\begin{enumerate}
  \item[e')] $\mathrm{Ext}^n(A/\mathfrak{N}, A)$ est comonogène
  \item[e'')] $A/(x_1 \ldots x_n)A$ est comonogène [i.e.\
    \uncertain{autodual}]
\end{enumerate}

\textbf{Corollaire 2.} Supposons $A$ quotient d'un régulier $A'$. Alors
ces conditions équivalent à :
\[
  \mathrm{Ext}^{n'-n}_{A'}(A, A') \simeq A
\]

\textbf{Remarque.} \struck{\ill{}} Il se peut que dans \uncertain{toutes}
les conditions, l'hyp.\ Cohen-Macaulay soit surabondante.

\textbf{Déf 2.} Anneau de Gorenstein $=$ Anneau Coh.\ Mac.\ $A$ tel que
\uncertain{$A$} \struck{\ill{}} \uncertain{fond.}

Ex : les réguliers.

\textbf{Proposition 17.} (i) Un localisé d'un Gorenstein
(\uncertain{supposé} quotient d'un régulier) est Gorenstein.

(ii) Kif kif pour un quotient par l'idéal engendré par une partie d'un
syst.\ de paramètres.

\textbf{Cor.} Le quotient d'un régulier par un idéal engendré par une
\uncertain{partie d'un} syst.\ de paramètres.
\note{« partie d'un » est ajouté au-dessus de la ligne ; la phrase
s'arrête là, sans verbe.}

\section{[Récapitulation des conditions ; le théorème]}
\note{titre de l'éditeur, entre crochets.}

\page{35}
\note{feuillet de récapitulation, d'une écriture plus posée : les
données fixées, puis cinq conditions (i) à (v) sur $M$, reliées à gauche
par des doubles flèches $\Downarrow$ de (i) à (ii) et de (i) à (iv), et par
une double flèche pointillée remontant de (i bis) à la réserve « [si $M$
équidim.] ».}

$M$ module de type fini sur l'anneau local $A$ de dim $n$

$k$ corps résiduel de $A$

$\mathfrak{N}$ idéal de définition fixé

$x_1 \ldots x_n$ syst.\ de paramètres fixé

$\mathfrak{p}$ un idéal premier de $A$ de corang $n$, fixé

$N$ variable de long.\ finie, $N_0 \neq 0$ fixé

\struck{\ill{}} (i) $\mathrm{Ext}^{n-1}_A(k, M) = 0$ et l'une des
conditions suivantes a) à d) [qui sont équivalentes moyennant la
première, parce qu'on a $\mathrm{Ext}^n(N, M) \simeq
\mathrm{Hom}(N, H^n(M))$ etc.]
\begin{enumerate}
  \item[a)] $\mathrm{Ext}^{n+1}(k, M) = 0$, \quad
    $\mathrm{Ext}^n(k, M) \simeq k$
  \item[$*$b)] $N \mapsto \mathrm{Ext}^n(N, M)$ est dualisant
  \item[$*$c)] $N \mapsto \mathrm{Ext}^n(N, M)$ est exact,
    $0 < \mathrm{lg}_{\mathfrak{p}} M \leq \mathrm{lg}_{\mathfrak{p}} A$
  \item[b')] $H^n(M)$ dualisant
  \item[c')] $H^n(M)$ injectif,
    $0 < \mathrm{lg}_{\mathfrak{p}} M \leq \mathrm{lg}_{\mathfrak{p}} A$
  \item[d)] $\mathrm{Ext}^{n+1}(k, M) = 0$,
    $0 < \mathrm{lg}_{\mathfrak{p}} M \leq \mathrm{lg}_{\mathfrak{p}} A$
\end{enumerate}
\note{les lettres b'), c'), d) sont repassées sur de premières lettres ;
une lettre noircie sous a) ne se lit pas.}

(i bis) $\mathrm{Ext}^{n-1}_A(k, M) = 0$ et l'une des \struck{deux}
cond.\ équiv.\ suivantes
\begin{enumerate}
  \item[$*$e)] $\mathrm{Ext}^n(k, M) \simeq k$, $M$ fidèle
  \item[\struck{$*$f)}] \struck{$\mathrm{Ext}^n(A/\mathfrak{N}, M)$ est
    comonogène, $M$ fidèle}
\end{enumerate}
\marginal{[si $M$ \struck{\ill{}} équidimen.]}
\note{la ligne f) est barrée de deux traits ; la réserve entre crochets
est écrite au-dessus de (i bis), « équidimen. » au-dessus d'un mot
noirci.}

(ii) Sous l'une des formes suivantes, qui sont équivalentes
\begin{enumerate}
  \item[a)] $N \mapsto \mathrm{Ext}^n(N, M)$ est dualisant
  \item[b)] $N \mapsto \mathrm{Ext}^n(N, M)$ exact,
    $0 < \mathrm{lg}_{\mathfrak{p}} M \leq \mathrm{lg}_{\mathfrak{p}} A$
\end{enumerate}

(iii) $\mathrm{Ext}^n(k, M) \simeq k$, $M$ fidèle
\struck{\ill{}}

(iv) $\mathrm{Ext}^n(A/\mathfrak{N}, M) \simeq D(A/\mathfrak{N})$,
$M$ fidèle \struck{\ill{}}

peut-être suffit-il ici que $\mathrm{Ext}^n(A/\mathfrak{N}, M)$ soit
comonogène et $M$ de support $\mathrm{Spec}(A)$

\marginal{N.B. \uncertain{Ici} si $M$ est Coh.\ Mac.\ \struck{\ill{}}
\ill{}, \ill{} $\mathrm{Ext}^n(N, M) =
\mathrm{Hom}(N, H^n(M))$ \ill{} \ill{} \ill{}, \ill{} \ill{} \ill{}
\ill{} \ill{}}
\note{note à droite de (iii) et (iv), qu'une accolade réunit ; « $A$ et »
est ajouté au-dessus de « $M$ » ; une ligne est biffée de traits ondulés,
et deux traits ondulés prolongent (iii) et (iv) jusqu'à elle.}

(v) L'annulateur de \struck{$\mathfrak{N}$} dans $M/x_1M + \cdots + x_nM$
est de dim 1, $M$ est fidèle

Ces conditions sont vraies si
\[
  \boxed{A \ \text{est Coh.\ Mac.\ et} \ M \simeq \Omega_A}
\]

Pour prouver la réciproque, il suffirait de montrer que dans les cas
\struck{(i)} (i bis) à (iii) (si on prend (i bis) \uncertain{faible}) que
\uncertain{si $n > 0$} l'annulateur de $\mathfrak{m}$ dans $M$ est nul,
\struck{dans le cas (iv), qu'il en est de même avec $M$ et $A$}. Dans (v),
\struck{\ill{} \ill{} si $n > 0$ et si} commençons par prouver que si
$n > 0$, et \struck{la première} l'\uncertain{hypothèse} condition de (v),
est vérifiée, alors l'annulateur de $\mathfrak{m}$ dans $M$ est nul [d'où
déjà que $M$ est Coh.-Mac.\ par récurrence, puis que
$\mathrm{Ext}^n(A/\underline{x}, M) \simeq M/\underline{x}M$ et
$\mathrm{Ext}^n(k, M) \simeq k$, d'où le fait que $H^n(M)$ est le dual d'un
$A/\mathfrak{I}$, $\mathfrak{I}$ intersection d'idéaux primaires de
corang $n$
\note{« si $n > 0$ » est ajouté au-dessus de la ligne ; la phrase
continue au feuillet suivant. La lettre notée $\mathfrak{m}$ est un petit
signe rond, peut-être l'idéal maximal.}

\page{36}
\uncertain{(remplaçant $A$ par $A/\mathfrak{I}$)}], d'où
\uncertain{(ii)} $+$ le fait que $M$ est Coh.\ Mac. Or de ceci on déduit
facilement que $A$ est Coh.\ Mac.\ et $M \simeq \Omega_A$. \uncertain{On}
peut essayer une voie analogue pour (iv), en y \uncertain{divisant} de la
précédente hyp.\ [que l'annulateur de $\mathfrak{m}$ \ill{} \ill{} $M$
\emph{et} $A$ \struck{sont} \uncertain{nuls}] \uncertain{$+$ la
condition \ill{} \ill{}}, d'où par récurrence que $M$ \emph{et} $A$ sont
Coh.\ Mac., \struck{\ill{}} \uncertain{ce} qui ramène comme
précédemment à un cas de (i)
\note{« (remplaçant $A$ par $A/\mathfrak{I}$) » est écrit au-dessus de la
première ligne et relié par un crochet à ce qui suit « d'où » ; « $+$ la
condition … » est ajouté au-dessus de « hyp. ».}

\uncertain{Démonstration} des équivalences dans (i)

b $\Rightarrow$ b' \quad c $\Rightarrow$ c', \quad a $\Rightarrow$ b
\struck{\ill{}} triviaux

b $\Rightarrow$ c car b $\Leftrightarrow$ b', donc
$D(H^n(M)) \simeq \hat{A}$ et
$\mathrm{lg}_{\mathfrak{p}}(A) = \mathrm{lg}_{\mathfrak{p}}(M)$

c' $\Rightarrow$ b' car c' $\simeq$ \ill{}
$D(H^n(M)) \simeq \hat{A}^m$ et
$m\, \mathrm{lg}_{\mathfrak{p}}(A) = \mathrm{lg}_{\mathfrak{p}}(M)$ d'où
$m = 1$

c $\Rightarrow$ d, car $N \mapsto \mathrm{Ext}^{n+1}(N, M)$ est exact à
gauche, d'où $\mathrm{Ext}^{n+1}(k, M) \to H^{n+1}(M)$ inj.
\marginal{\struck{b'} \ill{}}

d $\Rightarrow$ c' trivial, \struck{\ill{}} // donc
\qquad a $\Rightarrow$ b $\Leftrightarrow$ b' $\Leftrightarrow$ c
$\Leftrightarrow$ c' $\Leftrightarrow$ d

b $\Rightarrow$ $\mathrm{Ext}^n(k, M) \simeq k$,
$\mathrm{Ext}^n(A/\mathfrak{N}, M) \simeq D(A/\mathfrak{N})$ et
b') $\Rightarrow$ $H^n(M)$ est fidèle donc $M$ est fidèle, d'où
(b, b') $\Rightarrow$ e et f, \quad \uncertain{donc} (b, b', d)
$\Rightarrow$ a

Reste à prouver que f) $\Rightarrow$ a). Or \uncertain{si $M$ est Coh.\
Mac.}
\[
  \begin{array}{l}
    \mathrm{Ext}^n(A/\mathfrak{N}, M) = \mathrm{Hom}(A/\mathfrak{N}, M) \\
    \mathrm{Ext}^n(k, M) = \mathrm{Hom}(k, M)
  \end{array}
\]
\note{les deux $\mathrm{Hom}$ sont écrits tels quels ; on attend
l'argument $M/\underline{x}M$ ou analogue, que la page ne porte pas.}
d'où $\mathrm{Ext}^n(k, M) =$ annulateur de $\mathfrak{m}$ dans
$\mathrm{Ext}^n(A/\mathfrak{N}, M)$ donc f) $\Rightarrow$ e), et e)
$\Rightarrow$ $H^n(M) \subset D(A)$. \struck{Donc \ill{}} \ill{}
\uncertain{l'annulateur} \ill{}

\struck{Soit $\mathfrak{I}$ l'annulateur de $H^n(M)$, on a
$H^n(M) \simeq D(A')$, $A' = A/\mathfrak{I}$. Or $\mathfrak{I}$ est
l'intersection des idéaux primaires de corang $n$ \uncertain{associés à}
$M$. Et $H^n(M)$ \ill{} \ill{} \ill{} \uncertain{si on remplace} $M$ par
\ill{} $M/M_1$, \ill{} $M_1$ est l'intersection \ill{} \ill{} \ill{} \ill{}
$M_{\mathfrak{p}}$ ($\mathfrak{p}$ de corang $n$).}
\note{ce passage est encadré et biffé de traits obliques et horizontaux.}

\uncertain{et} $H^n(M)$ est celui de $M$ ($M$ Coh.\ Mac.\ donc
\emph{équidimensionnel}) donc $M$ étant fidèle, il est nul ; \uncertain{ce}
\uncertain{qui} signifie (supposons $A$ complet, ce qui est loisible) que
$H^n(M) = D(A)$, i.e.\ b').

\struck{Donc \ill{} \ill{}} Il reste donc essentiellement les

\textbf{Problème A.} Si dans $M/x_1M + \cdots + x_nM$, l'annulateur de
$\mathfrak{m}$ est de dim 1, et [\struck{dim} $M = n$] prouver que dans
$M$ il est nul

\textbf{Problème B.} Si dans $\mathrm{Ext}^n(A/\mathfrak{N}, M)$, \ill{}
annulateur de $\mathfrak{m}$ \uncertain{est} \ill{} de dim 1, et
[$\uncertain{\dim} M = \dim A = \ill{}$] prouver que dans $M$
\uncertain{et $A$} il est nul
\note{« dans » est écrit au-dessus de la ligne dans les deux problèmes, et
les crochets de fin de ligne, en marge droite, sont en partie hors du
feuillet ; la dernière ligne est coupée par le bord.}

\page{37}
On en déduirait des \emph{corollaires} : dans les deux cas, $M$ est
Coh.-Mac., et si on désigne par $A'$ son anneau d'homothéties, $A'$ est
Coh.\ Mac., et $M \simeq \Omega_{A'}$.

\textbf{Problème C.} Supposons
$\mathrm{Ext}^{n-1}(k, M) = \mathrm{Ext}^{n+1}(k, M) = 0$, (la dernière
condition équivalant à : $H^n(M)$ injectif) \uncertain{ou} :
$N \mapsto \mathrm{Ext}^n(N, M)$ exact). Est-il vrai que $A$ est Coh.\
Mac.\ et $M \simeq \Omega_A^m$\,?
\marginal{$\mathrm{Ext}^n(k, M) \neq 0$ i.e.\ $H^n(M) \neq 0$}
\note{l'hypothèse additionnelle est écrite au-dessus de la ligne, dans une
boucle qui la rattache à « Supposons » ; « ou : $N \mapsto \ldots$ exact »
est en interligne sous la parenthèse.}

[On notera qu'on a un hom.\ naturel $M \to \Omega_A^m$, il faut montrer
que c'est un isomorphisme]. \emph{\uncertain{Sinon} \uncertain{encore}},
on est ramené (par la dévissage : \uncertain{le} Serre) à prouver que
l'annulateur de $\mathfrak{m}$ dans $M$ \struck{est} \emph{et} $A$
\struck{est} \uncertain{sont} nuls.
\note{la page s'arrête là ; le reste du feuillet est blanc.}

\page{38}
\textbf{Théorème.} \uncertain{Les} \ill{} \uncertain{les} implications,
marquées sur le diagramme des conditions. De plus, si on remplace dans
(i) et (i bis) la condition $\mathrm{Ext}^{n-1}(k, M) = 0$ par :
$M$ est Coh.\ Mac., et si dans (v) on suppose de plus $M$ Coh.\ Mac.,
alors \struck{toutes} les conditions (i) (i bis) (v) deviennent
équivalentes à : $A$ est Coh.\ Mac.\ et $M \simeq \Omega_A$
(\uncertain{et} \uncertain{par} \uncertain{les} conditions ii iii iv, qui
pour $M$ Coh.\ Mac.\ sont des \uncertain{variantes} de (i) et (i bis))
\note{l'énoncé est encadré à gauche d'un trait ; le « diagramme des
conditions » est celui de la page 35.}

$M$ est équidim., donc (i) $\Leftrightarrow$ (i bis). Prouvons par
récurrence que alors $A$ est Coh.\ Mac.\ et $M \simeq \Omega_A$.
Pour $n = 0$ c'est trivial [$H^0(M) = M$ est dualisant].
Supposons $n > 0$, et le th.\ prouvé pour $n' < n$. \struck{\ill{}}
Comme $H^n(M)$ est fidèle, \uncertain{$M$ est fidèle, et \ill{}}
\struck{les idéaux premiers associés à $A$ sont de dim.\ $n$, donc et}
\struck{\ill{}} Coh.\ Mac., \uncertain{il} s'ensuit que les idéaux premiers
associés à $M$ et à $A$ sont les \uncertain{mêmes}, et qu'ils sont tous de
\struck{dimension} \uncertain{corang} $n$. Comme $n > 0$, il y a un $x$ qui
n'est contenu dans aucun d'eux, donc qui n'est pas div.\ de 0 dans $A$ ni
dans $M$. \struck{D'après} \uncertain{Ecrivons}, \uncertain{alors}
\[
  \mathrm{Ext}^{\bullet}_A(N, M) \Longleftarrow
  \mathrm{Ext}^{\bullet}_{A/x}\bigl(N, \mathrm{Ext}^{\bullet}_A(A/x, M)\bigr)
\]
Or $\mathrm{Ext}^0_A(A/x, M) = 0$,
$\mathrm{Ext}^1_A(A/x, M) = M/xM$, d'où
\[
  \mathrm{Ext}^i_A(N, M) \simeq \mathrm{Ext}^{i-1}_{A/x}(N, M/xM)
\]
\uncertain{pour} tout $N$. Par suite, la condition (i a)
\note{« corang » est écrit au-dessus de « dimension » biffé. La phrase
continue au feuillet suivant.}

\page{39}
reste vraie pour le module $M/xM$ sur $A/xA$. Or ce dernier est
\uncertain{encore} Coh.\ Mac., et l'hyp.\ de récurrence s'applique.
Donc \struck{$M/xM$} $A/xA$ est Coh.\ M., donc $A$ l'est. D'autre part
$M/xM$ \struck{$\simeq D_{A/xA}$ \ill{}} est fond.\ pour $A/xA$, d'où
\[
  \begin{array}{c}
    D(M/xM) = \text{annulateur de } x \text{ dans } D(M) \\
    \wr| \\
    H^n(A/xA) = \text{annulateur de } x \text{ dans } H^n(A) .
  \end{array}
\]
Donc l'hom.\ naturel [\uncertain{déduit} de
$H^n(M) \simeq H^n(A) \otimes M = D(A)$]
\[
  H^n(A) \longrightarrow D(M)
\]
est tel qu'il induit un isom.\ sur les \uncertain{éléments} annulés par
$x$, donc (\uncertain{comme} \ill{}) il est injectif, ou par
transposition,
\[
  M \longrightarrow \Omega_A
\]
est surjectif. D'autre part il induit un isom.\ sur les $H^n$, i.e.\ le
$H^n$ du noyau est \ill{}, donc ce noyau est de dim $< n$. Comme $M$ est
équidim., ce noyau est nul.
\note{les crochets de « déduit de $H^n(M) \simeq \ldots = D(A)$ » sont
écrits jusqu'au bord droit du feuillet.}

\struck{Prouvons maintenant (v), en supposant $M$ Coh.\ Mac.}

\struck{Posons} Supposons \struck{$\mathrm{Supp}\, M = \mathrm{Spec}(A)$},
\uncertain{que dans} $M/x_1M + \cdots + x_nM$ l'annulateur de $\mathfrak{m}$
\uncertain{soit} de dim 1, $M$ \struck{Coh.\ Mac.} Coh.\ Mac.\ fidèle.
\struck{Prouvons} Alors
($\mathrm{Ext}^n(k, M) = \struck{\mathrm{Hom}(A/\underline{x}, M/\underline{x}M)}$
$\mathrm{Hom}(k, M/\underline{x}M) \simeq k$) (cf.\ § 1, prop.\ 4,
cor.\ 1) donc on est ramené à (i et e)
\note{les biffures de cette fin de page sont reprises en interligne ; dans
$\mathrm{Ext}^n(k, M)$ le $k$ est écrit sur un « $A/\underline{x}$ »
noirci. La référence « § 1, prop. 4, cor. 1 » renvoie à un feuillet hors de
ce lot.}

\page{40}
\note{nouvelle récapitulation, écrite d'abord à l'encre noire puis reprise
à l'encre bleue et au crayon, avec de nombreuses surcharges, des flèches
et des notes obliques dans les deux marges ; on donne le corps de la page
et ce qui se lit des marges.}

$A$ anneau local, quotient d'un régulier $A'$. \quad $n = $ dim.\ Krull $A$

$\underline{k} = A/\mathfrak{m}(A)$

$\mathfrak{N}$ un idéal de définition

$(x_1, \ldots, x_n)$ un système de paramètres.

$\mathfrak{p}$ un idéal premier de dim $n$.

\emph{Conditions}
\begin{enumerate}
  \item[(i)] $\mathrm{Ext}^n(\underline{k}, M) \simeq k$,
    $\mathrm{Ext}^{n+1}(\underline{k}, M) = \mathrm{Ext}^{n-1}(\underline{k}, M)
    = 0$
  \item[(ii)] $\mathrm{Ext}^n(\underline{k}, M) \simeq k$ et $M$
    \emph{fidèle} [\ill{} \ill{} \ill{} $M$]
  \item[\struck{(iii)}] \struck{$\mathrm{Ext}^{n-1}(\underline{k}, M) =
    \mathrm{Ext}^{n+1}(\underline{k}, M) = 0$,
    $0 < \mathrm{lg}_{\mathfrak{p}}(M) \leq \mathrm{lg}_{\mathfrak{p}} A$}
  \item[\struck{(i bis)}] \struck{$\mathrm{Ext}^{n-1}(\underline{k}, M) = 0$,
    $M$ fidèle, $0 < \mathrm{lg}_{\mathfrak{p}}(M) \leq
    \mathrm{lg}_{\mathfrak{p}}(A)$}
  \item[\struck{(iv)}] \struck{$\mathrm{Ext}^{n-1}(\underline{k}, M) = 0$
    et $\mathrm{Ext}^n(A/\mathfrak{N}, M) \simeq A/\mathfrak{N}$}
  \item[(ii bis)] $\mathrm{Ext}^n(A/\mathfrak{N}, M) \simeq
    \widehat{A/\mathfrak{N}}$ et $M$ \emph{fidèle} [faut-il
    $\mathrm{Ext}^{n-1}(k, M) = 0$\,? ou $A$ équidim\,?]
    \uncertain{suffit-il} que $\mathrm{Ext}^n(A/\mathfrak{N}, M)$ soit
    \emph{comonogène}\,?
  \item[(iii)] $\mathrm{Ext}^{n-1}(\underline{k}, M) = 0$, et l'une des
    cond.\ équivalentes
    \[
      \left\{\begin{array}{l}
        \struck{\ill{}} \\
        \mathrm{Ext}^{n+1}(\underline{k}, M) = 0 \\
        H^n(M) \ \text{dualisant} \\
        H^n(M) \ \text{injectif}
      \end{array}\right.
    \]
    et $0 < \mathrm{lg}_{\mathfrak{p}}(M) \leq \mathrm{lg}_{\mathfrak{p}}(A)$
  \item[(iv)] La fonction $N \mapsto \mathrm{Ext}^n(N, M)$ est
    dualisant
  \item[(iv bis)] La fonction $N \mapsto \mathrm{Ext}^n(N, M)$ est exact,
    $0 < \mathrm{lg}_{\mathfrak{p}}(M) \leq \mathrm{lg}_{\mathfrak{p}}(A)$
    ($N$ de long.\ finie)
  \item[(v)] \uncertain{faible} L'annulateur de $\mathfrak{m}$ dans
    $M/\sum x_i M$ est de dim 1 et $M$ est fidèle
\end{enumerate}
\note{la condition (iii) et deux reprises, (i bis) en bleu et (iv) en
noir, sont encadrées ensemble et barrées de traits en zigzag ; les
lettres a, b, c cerclées en marge droite renvoient aux trois conditions
de l'accolade de (iii). Une grande accolade bleue, à gauche, embrasse (i)
à (ii bis).}

\marginal{$\mathrm{Ext}^n$ et $\mathrm{Ext}^{n+1}$ \ill{} \ill{} \ill{}
ou $\mathrm{Ext}^n$ et $\mathrm{Ext}^{n-1}$ ou $\mathrm{Ext}^{n-1}$ et
$\mathrm{Ext}^{n+1}$ \ill{} \ill{} $A/\mathfrak{p}$, $A$ et
$M_{\mathfrak{p}}$ \ill{}, $A/\mathfrak{I}$, $\Omega_A$, \ill{} \ill{}
\ill{}, $A^m$, \ill{} \ill{}}
\note{note de la marge droite, en regard de (i) et (ii), reliée par une
double flèche au crochet de (i) ; au-dessus, trois lignes brèves ne se
lisent pas.}

\[
  \boxed{A \ \text{est C.-M., et} \ M \ \text{est} \ \simeq \Omega_A}
\]

\textbf{Pb.} \quad $M \neq 0$,
$\mathrm{Ext}^{n-1}(k, M) = \left\{\begin{array}{l}
  \mathrm{Ext}^{n+1}(k, M) = 0 \\
  H^n(M) \ \text{injectif}
\end{array}\right.$
$\Longrightarrow$ $A$ est C.M.\ et $M \simeq \Omega_A^m$
\marginal{mais une condition \ill{} \ill{} \ill{} \ill{} \ill{}
$M \simeq \Omega_A \times A/\mathfrak{I}$}
\note{le problème est encadré à gauche d'un trait ; une flèche double
descend vers lui depuis une note oblique au crayon, « $M \neq 0$,
$\mathrm{Ext}^n(N, M)$ exact » ; la note sur $\Omega_A \times A/\mathfrak{I}$
est dans un encadré en coin, en bas à droite, avec un mot biffé.}

\marginal{\ill{} $H^n(M)$ \ill{} \ill{} \ill{} \ill{} \ill{} $M$ \ill{}
\ill{} $H^n(M)$ \ill{} \ill{} $H^n(M)$ et \ill{} $M$ équidim.\ \ill{}
$\mathrm{Ext}^{n-1}(k, M)$ \ill{} (vi) \ill{} \ill{} $\Omega_{A/\mathfrak{I}}$
\ill{}}
\note{notes obliques de la marge gauche, en deux couches (bleu, puis
noir), autour d'un « (vi) » encadré ; on n'en lit que des fragments.}

\struck{Le seul cas \ill{} \ill{} $\mathrm{Ext}^{n-1}(k, M) \neq 0$,
\uncertain{$M \neq 0$}, \ill{} \ill{} \ill{} $A/\mathfrak{I}$ (\ill{}
équidim.\ \uncertain{associé} à $A$) est C.M.\ et [\ill{} \ill{}
$\mathrm{Ext}^n(k, M) = k$] si $m = 1$ \ill{} \ill{} \ill{}
$\mathrm{Ext}^n(k, M) \simeq k$ est \uncertain{celui} \ill{}
$A/\mathfrak{I}$ est C.-M.\ et \ill{} $M = \Omega_{A/\mathfrak{N}}$, où
$\mathfrak{N}$ est un idéal de $A$ tel que $A/\mathfrak{N}$ soit C.-M.,
\ill{}, \ill{} $m = 0$, \ill{} \ill{} \ill{} $\hat{A}$}
\note{ce dernier paragraphe est barré de trois longs traits obliques ;
une flèche au crayon en part vers une ligne au pied de la page,
« \ill{} \ill{} \ill{} $M$ \ill{} », elle-même biffée.}

\end{document}
